\documentclass[11pt,letterpaper]{article}
\usepackage{technical-design-spec}
% Shared non-semantic helper definitions; one semantic module is loaded by each root.
\newcommand{\PTSource}[1]{\par\smallskip{\footnotesize\color{Meta}\textbf{Book crosswalk:} #1}\par\smallskip}
\newcommand{\PTOutcome}[1]{\begin{tipbox}[Competency outcome]#1\end{tipbox}}
\newcommand{\PTEvidence}[1]{\begin{notebox}[Required evidence]#1\end{notebox}}
\newcount\ptlinecount
\newcommand{\PTLines}[1]{\par\ptlinecount=0\loop\ifnum\ptlinecount<#1\par\noindent\rule{\linewidth}{0.25pt}\vspace{0.22cm}\advance\ptlinecount by1\repeat\par}
\newcommand{\PTField}[1]{\noindent\textbf{#1}\quad\rule{0.66\linewidth}{0.3pt}\par\vspace{0.22cm}}
\newcommand{\PTRating}{\noindent\textbf{Rating:}\quad 0\quad 1\quad 2\quad 3\quad 4\qquad\textbf{Evidence ref:}\rule{0.30\linewidth}{0.3pt}\par\vspace{0.14cm}}
\newcommand{\PTTask}[2]{\par\medskip\noindent\textbf{Task #1.} #2\par}

\setDocTitle{Formal Grammars and Automata: Field Guide}
\setDocSubtitle{Derivations, equivalence, FIRST/FOLLOW, DFA construction and fixed points}
\setDocVersion{1.0}\setDocStatus{Engineering education / source-aligned}
\setDocOwner{Compiler Engineering}
\setDocDate{30 October 2026}
\setDocSummary{Engineering reference and source crosswalk.}
\begin{document}\makedoctitle\makedocsummary
\tableofcontents\clearpage

\section{Derivation and notation}
Use $\Rightarrow$ for one production application, $\Rightarrow^*$ for zero or more, and $\epsilon$ for the empty string. For $S\to aSb\mid\epsilon$, $S\Rightarrow aSb\Rightarrow aaSbb\Rightarrow aabb$. The generated language is not regular because matching unbounded counts requires unbounded memory. Its grammar is context-free.\PTSource{Chapters 2--3.}
\section{Grammatical ambiguity}
A grammar is ambiguous if one terminal string has at least two distinct parse trees. The grammar $E\to E+E\mid E*E\mid n$ is ambiguous for $n+n*n$. Rewriting the grammar into precedence layers may remove this ambiguity but does not prove that all equivalent grammars are unambiguous.\PTSource{Chapter 3.}
\section{Grammar sanitation}
Compute reachability from $S$ by graph traversal. Compute productivity by a fixed point: a nonterminal is productive if it has a production whose right-hand symbols are terminals or already productive nonterminals. A nonterminal may be reachable but never produce a terminal word. Record cycles and nullable productions separately.\PTSource{Chapter 2.}
\section{Automata construction}
A DFA is $(Q,\Sigma,\delta,q_0,F)$ with total $\delta:Q\times\Sigma\to Q$; model dead states explicitly when an input alphabet is complete. For NFA conversion, use subset construction with epsilon closures; only reachable subsets need be materialized. Minimize a DFA by partition refinement, preserving distinctions between accepting and nonaccepting states.\PTSource{Chapter 5.}
\section{FIRST and FOLLOW worked procedure}
Start with terminals: $\mathrm{FIRST}(a)=\{a\}$. Iterate production implications to a fixed point; for $A\to X_1\ldots X_k$, propagate the first nonnullable FIRST sets, and add $\epsilon$ only if all $X_i$ are nullable. Initialize $\$\in\mathrm{FOLLOW}(S)$, add $\mathrm{FIRST}(\beta)\setminus\{\epsilon\}$ to FOLLOW of a nonterminal before suffix $\beta$, and propagate FOLLOW of the left-hand side if $\beta$ is nullable.\PTSource{Chapter 8.}
\section{Normal-form transformation cautions}
CNF makes binary productions convenient for CYK, but nullable production elimination, unit-rule elimination and binarization may introduce helper nonterminals and alter original parse trees. Preserve transformation provenance if downstream semantic actions need the original grammar structure.\PTSource{Chapter 4.}

\end{document}
